\section{Voucher Validation}
Voucher validation includes status, validity window, optional concert/category
scope, remaining global quota, and one-use-per-user policy. A global voucher
has NULL scope columns; a scoped voucher must match the booking concert and at
least one requested category. The server calculates the discount; the client
cannot submit an authoritative discount amount.

\section{Atomic Global Quota Update}
\begin{lstlisting}[language=SQL,caption={Atomic voucher quota reservation}]
UPDATE vouchers
SET used_count = used_count + 1
WHERE id = :voucher_id
  AND status = 'ACTIVE'
  AND starts_at <= NOW()
  AND expires_at > NOW()
  AND used_count < usage_limit
RETURNING used_count;
\end{lstlisting}

\section{Per-User Anti-Abuse Rule}
For a one-use-per-user policy, a partial unique index on
\code{(voucher\_id, user\_id)} where redemption status is not RELEASED provides
a database-backed guarantee. A retry of the same logical booking does not
create another redemption because idempotency is applied before the booking can
be duplicated.

\section{Release Semantics}
Voucher quota is reserved at booking creation, represented by a
\code{RESERVED} redemption, and counted in \code{used\_count}. Confirmation
transitions it to \code{CONSUMED}. Expiration/cancellation transitions only a
\code{RESERVED} redemption to \code{RELEASED}, decrements the quota in the same
transaction, and keeps the audit row. Repeating release affects zero rows, so
the allocation cannot be decremented twice.
